\chapter{Itchy Anus (Pruritus Ani)}
\index{Itchy Anus (Pruritus Ani)}

\section*{Definition and Overview}
An itchy anus, medically called pruritus ani, is persistent itching or burning around the back passage. It is a common and distressing symptom that can disrupt sleep and daily life, and people often delay seeking help because of embarrassment. Most cases are caused by minor skin irritation, moisture, or hygiene problems, and respond well to simple measures. Some cases have a specific cause, such as haemorrhoids, anal fissures, skin conditions, threadworms, or fungal infection. Persistent itching with bleeding, pain, or a change in bowel habit should always be assessed, because serious conditions need to be excluded.

\section*{Epidemiology}
An itchy anus affects around 1 to 5 per cent of people at some time and is more common in men, in people with haemorrhoids or bowel conditions, and in those who sweat heavily or are overweight. It peaks in middle age. Threadworms are a common cause in children and can spread through families. Many people self-treat or cope without seeking help, and it is one of the most common reasons for referral to colorectal specialists. With correct diagnosis and management, most cases resolve.

\section*{Aetiology and Pathophysiology}
The skin around the anus is delicate and prone to irritation. The most common causes are excessive moisture from sweating or poor drying, and faecal soiling from incomplete wiping, which irritates the skin. Over-cleaning with soaps, wipes, and vigorous rubbing strips the skin's protective barrier and worsens itching, creating a scratch-itch cycle. Specific causes include haemorrhoids, which leak mucus and cause soiling; anal fissures; skin conditions such as eczema and psoriasis; fungal infections (tinea); and threadworms, which lay eggs around the anus at night, causing intense nocturnal itching. Less commonly, the itch is caused by contact allergy, diabetes, or, rarely, more serious disease, so persistent symptoms should be investigated.

\section*{Clinical Features}
\begin{itemize}
    \item Itching, burning, or soreness around the back passage
    \item Itching that is worse at night, which may suggest threadworms
    \item A feeling of moisture or soiling
    \item Redness, thickening, or cracking of the surrounding skin
    \item A lump at the anus, which may indicate a haemorrhoid or skin tag
    \item Bleeding with wiping, suggesting a fissure or haemorrhoid
    \item Disturbed sleep and difficulty concentrating from persistent itch
    \item Symptoms that flare after sweating, exercise, or certain foods and drinks
\end{itemize}

\section*{Differential Diagnosis}
The cause of anal itching must be identified because treatments differ.
\begin{itemize}
    \item \textbf{Idiopathic pruritus ani:} no cause found, usually related to hygiene and moisture
    \item \textbf{Haemorrhoids:} internal or external piles with mucus discharge
    \item \textbf{Anal fissure:} a painful tear in the anal lining
    \item \textbf{Threadworms and pinworms:} especially in children, with night-time itching
    \item \textbf{Skin conditions:} eczema, psoriasis, and contact dermatitis
    \item \textbf{Fungal infection:} tinea of the perianal skin
    \item \textbf{Diabetes:} can cause generalised and local itching
    \item \textbf{Rarely, more serious disease:} such as anal cancer, warranting examination
\end{itemize}

\section*{When to Seek Medical Care}
Medical assessment is appropriate for itching that is persistent, severe, or interfering with sleep; for itching with bleeding, pain, a lump, discharge, or a change in bowel habit; and for symptoms in children that suggest threadworms, as the whole family may need treatment. People with diabetes or weakened immunity should seek earlier review.

\textbf{Contact your local emergency services for:}
\begin{itemize}
    \item Severe pain or a painful, tense swelling at the anus, which may indicate a thrombosed haemorrhoid or abscess
    \item Heavy rectal bleeding with faintness or dizziness
    \item Fever with severe pain around the anus
    \item Inability to pass urine or stool
\end{itemize}

\section*{Diagnosis and Investigations}
\begin{itemize}
    \item \textbf{History and examination:} inspection of the perianal skin and anus
    \item \textbf{Digital rectal examination:} to feel for haemorrhoids, fissures, and masses
    \item \textbf{Proctoscopy:} direct inspection of the anal canal
    \item \textbf{Stool testing:} for threadworms, using the tape test or stool sample
    \item \textbf{Skin scrapings:} for fungal infection where suspected
    \item \textbf{Blood tests:} for diabetes or other conditions where indicated
    \item \textbf{Colonoscopy:} for persistent bleeding or other warning features
\end{itemize}

\section*{Management}
\begin{itemize}
    \item \textbf{Hygiene measures:} gentle cleansing with water, thorough drying, and avoiding harsh soaps and wipes
    \item \textbf{Moisture control:} cotton underwear, avoiding tight clothing, and using barrier creams
    \item \textbf{Breaking the scratch-itch cycle:} short courses of soothing or steroid creams under medical advice
    \item \textbf{Treatment of the cause:} creams or banding for haemorrhoids, treatment for fissures and threadworms
    \item \textbf{Antifungal or antibacterial treatment:} for infection
    \item \textbf{Dietary review:} reducing spicy foods, caffeine, and alcohol if they trigger symptoms
    \item \textbf{Managing constipation:} fibre, fluids, and avoiding straining
    \item \textbf{Specialist referral:} for persistent or unexplained symptoms
\end{itemize}

\section*{Complications}
\begin{itemize}
    \item Scratch-itch cycle with worsening skin damage
    \item Lichenification: thickened, leathery skin from chronic scratching
    \item Secondary bacterial or fungal infection
    \item Sleep disturbance and reduced quality of life
    \item Embarrassment and delay in seeking care
    \item Missed underlying conditions when symptoms are dismissed
\end{itemize}

\section*{Prognosis and Outcomes}
Most cases of anal itching resolve with simple hygiene measures and treatment of the underlying cause within days to weeks. Breaking the scratch-itch cycle is the key, and this usually responds to gentle care and short-term treatment. Threadworms and fungal infections are readily cured. Persistent or unexplained itching warrants specialist review, which identifies the cause in most people and leads to effective treatment.

\section*{Prevention and Self-Care}
\begin{itemize}
    \item Clean gently with water rather than scrubbing with soap
    \item Dry the area thoroughly after washing and bathing
    \item Avoid scented wipes, soaps, and vigorous rubbing
    \item Wear loose cotton underwear and avoid tight clothing
    \item Use barrier creams if the skin is moist or sore
    \item Do not scratch, and keep nails short
    \item Treat constipation and avoid straining
    \item Reduce spicy foods, caffeine, and alcohol if they worsen symptoms
    \item Seek medical advice for persistent symptoms or bleeding
    \item If threadworms are found in one family member, treat the whole family
\end{itemize}

\section*{Sources and Further Reading}
The following reputable sources provide further information for healthcare students and patients seeking reliable, current advice about an itchy anus.
\begin{itemize}
    \item Healthdirect Australia. \textit{Itchy anus (pruritus ani).} https://www.healthdirect.gov.au/
    \item Gastroenterological Society of Australia. \textit{Anorectal conditions.} https://www.gesa.org.au/
    \item Better Health Channel. \textit{Anal itch.} https://www.betterhealth.vic.gov.au/
    \item NHS. \textit{Itchy bottom.} https://www.nhs.uk/conditions/itchy-bottom/
    \item Mayo Clinic. \textit{Anal itching.} https://www.mayoclinic.org/
    \item DermNet NZ. \textit{Pruritus ani.} https://dermnetnz.org/
\end{itemize}
